Wind and solar output, grid operation and electricity prices all depend on weather forecasts, but a forecast means something different to a trader, a commissioning engineer and a lender. This paper maps what a good weather forecast must deliver for each renewable technology, for each department across the asset life cycle (development, financing, engineering, construction, commissioning, operations, performance analysis, trading, grid operation) and for each horizon from seconds to decades. We collect the physical and statistical foundations in one notation and prove four results that govern forecast design: the power curve amplifies relative wind errors by a factor of at least three at partial load; site errors average out only down to a floor set by their mean correlation; averaging two models helps if and only if $k^2+2\rho k<3$ for their RMSE ratio $k$ and error correlation $\rho$; and under asymmetric imbalance costs the optimal bid is a quantile, not the mean. We review weather-driven price formation, operational and machine-learning weather models, simulations, free data and their licences, and the data-handling errors that most often break forecasts. A case study verifies \valNumRuns{} days of 00\,UTC forecasts from ECMWF IFS, ECMWF's machine-learning model AIFS and NOAA GFS against \valNumStations{} German stations (May to July 2025). At day-ahead lead times AIFS matches IFS for 10\,m wind (RMSE \valWindAifsRmseDayAhead{} vs.\ \valWindIfsRmseDayAhead\,m/s) and beats it for temperature, while GFS trails both. A multi-model mean with a running 30-day bias correction lowers the wind RMSE to \valWindMmmBcRmseDayAhead\,m/s and extends positive skill against climatology from \valWindIfsSkillZeroDays{} to \valWindMmmBcSkillZeroDays{} days. All four propositions hold on the data: the median partial-load error elasticity is \valElasticityMedian, averaging \valNumStations{} stations removes \valAggRelRmseAllPct\,\% of the error as predicted, and bidding the predicted quantile cuts the expected asymmetric cost by \valCostSavingPct\,\% against the point forecast. All data, code and figure tables are published with the paper.
